\begin{textsample}{POS Dim 4 – 2013\_c6 – Score 9.00 – 2013\_c6/enade\_2013\_enf\_7b93b4bd5c.txt}  \label{ex:f4_pos_041}
CAPÍTULO 3 ANÁLISE TÉCNICA DA PROVA Este capítulo tem por objetivo apresentar o desempenho dos estudantes concluintes de Enfermagem no ENADE/2013.
Para isso, foram calculadas as estatísticas básicas da prova em seu todo ( seção 3.1.1 ), bem como as estatísticas dos componentes relacionadas ao Componente de Formação Geral ( seção 3.1.2 ) e ao de Conhecimento Específico da Área ( seção 3.1.3 ).
Dadas as suas características, foram analisadas em separado, as questões objetivas ( seção 3.2 ) e as questões discursivas ( seção 3.3 ).
Para as questões objetivas, foram disponibilizados os índices de facilidade e o de discriminação bisserial, também em separado para os Componentes de Formação Geral ( seção 3.2.1 ) e de Conhecimento Específico ( seção 3.2.2 ).
De cada componente, uma das questões foi escolhida para exemplificar a análise gráfica relacionando as alternativas escolhidas pelos estudantes ( inclusive o gabarito ) com o número de acertos no componente.
O Anexo I apresenta a íntegra da análise gráfica para todas as questões objetivas.
Para cada uma das questões discursivas, os conteúdos dos tipos mais comuns de respostas dos alunos são apresentados e comparados com o padrão de resposta esperado ( ver Anexo VIII com o padrão de respostas ).
Tomando como base as duas questões discursivas do Componente de Formação Geral, a seção 3.3.1.6 apresenta comentários sobre a correção das respostas com respeito à Língua Portuguesa.
Nas tabelas, são apresentados o tamanho da população inscrita e de presentes, e as seguintes estatísticas das notas11 : média do desempenho na prova, erro padrão da média, desvio padrão, nota mínima, mediana e nota máxima.
Tais estatísticas contemplam o total de estudantes concluintes da Área de Enfermagem inscritos e presentes à prova do ENADE/2013, tendo -se em vista as seguintes agregações : ( a ) as Grandes Regiões e o país como um todo ; ( b ) a Categoria Administrativa ; e ( c ) a Organização Acadêmica.
Em relação aos gráficos de distribuição de notas, o intervalo considerado foi de 10 unidades, aberto à esquerda e fechado à direita, com exceção do primeiro intervalo, [ 0 ; 10 ], fechado em ambos os extremos.
Para os gráficos de distribuição das notas das questões discursivas, foram consideradas mais duas categorias : questão em branco e nota zero. 11 Essas estatísticas e outras estão definidas no Capítulo 1. 3.1 ESTATÍSTICAS BÁSICAS DA PROVA 3.1.1 Estatísticas Básicas Gerais A Tabela 3.1 apresenta as estatísticas básicas da prova por grande Região.
A população total de inscritos foi de 30.289.
Destes, 26.729 estiveram presentes, sendo 11,8 \% o índice de não comparecimento.
A Região de maior abstenção foi a Sudeste ( 14,9 \% ) e a de menor abstenção foi a Sul ( 4,2 \% ).
A média das notas da prova como um todo ( nas seções seguintes serão analisados os Componentes de Formação Geral e de Conhecimento Específico ) foi 51,2, sendo que os alunos da região Norte obtiveram a média mais baixa ( 48,7 ) e os da região Sul obtiveram a média mais alta ( 54,8 ).
As demais médias foram : 50,9 na região Nordeste, 51,2 na região Sudeste e 50,1 na região Centro-Oeste.
O desvio padrão para o Brasil como um todo foi 13,9, sendo o maior desvio padrão encontrado na região Nordeste ( 14,3 ) e o menor na região Sul ( 12,5 ), indicando uma menor dispersão das notas desta última região.
A região que obteve a maior nota máxima foi a Nordeste ( 88,3 ), ao passo que a região que atingiu a menor nota máxima foi a Norte ( 82,6 ).
A mediana do Brasil como um todo foi 51,8, sendo a maior mediana obtida na região Sul ( 55,9 ) e a menor obtida na Norte ( 49,1 ).
A nota mínima foi zero em todas as regiões, exceto na região Norte ( 5,8 ).
O comportamento das notas dos estudantes de todo o Brasil pode ser observado no Gráfico 3.1 que apresenta um histograma com a distribuição das mesmas.
Essa é uma distribuição unimodal com moda no intervalo ( 50 ; 60 ].
Destaca -se, também, o intervalo ( 40 ; 50 ] com frequência relativa próxima à modal.
O coeficiente de assimetria da distribuição das notas é negativo ( -0,27 ) e não podemos considerar que a distribuição é aproximadamente simétrica.
As distribuições por Grande Região também apresentam assimetria negativa, concentração pouco maior do lado direito do histograma e mais espalhada do lado esquerdo.
Os Gráficos 3.2, 3.3 e 3.4 apresentam informações referentes à média da nota final dos participantes, desagregando os resultados de acordo com, respectivamente, as Grandes Regiões do país, a Categoria Administrativa e a Organização Acadêmica.
Os gráficos apresentam o valor da média das notas como uma barra e os extremos do intervalo de confiança de 95 \% como linhas verticais unidas por uma linha horizontal na forma da letra H maiúscula.
Considerando -se o gráfico de notas segundo Grande Região ( Gráfico 3.2 ), observa -se que existe diferença estatisticamente significativa ao nível de 95 \% entre a maior média, obtida na região Sul ( 54,8 ) e a menor, obtida na região Norte ( 48,7 ).
Levando -se em conta os estudantes por Categorias Administrativas ( Gráfico 3.3 ), observa -se que existe diferença estatisticamente significante entre as médias das notas dos alunos de IES Públicas e Privadas.
Pode ser observado ainda que a média dos alunos de IES Públicas ( 59,2 ) é maior que a dos alunos de IES Privadas ( 49,5 ).
A diferença entre as médias das regiões Sul e Norte ( 6,1 ), a maior e a menor média regional, é inferior a diferença entre IES Públicas e Privadas ( 9,7 ).
Tendo como foco o Gráfico 3.4, que apresenta as notas médias das provas segundo Organização Acadêmica, constata -se que a diferença da média das notas dos estudantes provenientes de Universidades é estatisticamente significativa ao nível de 95 \% das médias das notas dos estudantes provenientes de Centros Universitários e Faculdades, assim como a diferença entre a média dos estudantes de Centros Universitários e Faculdades.
A maior média foi obtida pelos estudantes de Universidades ( 54,7 ), e a menor, pelos de Faculdades ( 48,5 ). 3.1.2 Estatísticas Básicas no Componente de Formação Geral A Tabela 3.2 apresenta as estatísticas básicas em relação ao componente da prova que avalia a Formação Geral dos estudantes concluintes.
Os alunos de todo Brasil obtiveram desempenho médio de 44,6.
Quanto à variabilidade, o desvio padrão das notas dos estudantes do Brasil como um todo foi 14,7.
A maior média foi obtida na região Sul ( 45,3 ), e a menor, na região Norte ( 42,8 ).
As demais médias foram : 44,7 nas regiões Nordeste e Sudeste e 44,0 na região Centro-Oeste.
Já o maior desvio padrão foi obtido na região Nordeste ( 15,1 ) e o menor na região Sul ( 14,0 ).
Os demais desvios padrões foram : 14,1 na região Norte, e 14,7 nas regiões Sudeste e Centro-Oeste.
A maior nota no Componente de Formação Geral da prova do ENADE foi obtida por pelo menos um aluno da região Sudeste ( 92,0 ) enquanto que a menor nota máxima foi obtida na região Norte ( 84,3 ).
Nas outras regiões as notas máximas foram : 88,8 na região Nordeste, 88,2 para a região Sul e 89,0 para a região Centro-Oeste.
A mediana do Brasil como um todo foi 45,0, sendo a menor mediana encontrada na região Norte ( 42,9 ) e a maior O Gráfico 3.5 propicia a avaliação do desempenho dos estudantes no Componente de Formação Geral a partir do histograma da distribuição das notas correspondentes.
A distribuição é unimodal, com moda em ( 40 ; 50 ], enquanto na prova como um todo a moda foi alcançada no intervalo ( 50 ; 60 ].
Nota -se, ainda, que no Gráfico 3.5 as notas apresentam uma maior dispersão do que no Gráfico 3.1 ( distribuição das notas da prova ), confirmado pela comparação dos desvios padrões : 13,9 para a nota da prova como um todo e 14,7 para o Componente de Formação Geral.
Para o Componente de Formação Geral, o coeficiente de assimetria da distribuição das notas dos estudantes é negativo ( -0,16 ).
Por ser um coeficiente pequeno, a distribuição apresenta uma leve concentração à direita e cauda maior à esquerda, porém podemos considerar que a distribuição é aproximadamente simétrica.
Em todas as Grandes Regiões os histogramas também possuem assimetria levemente negativa ( entre -0,20 para o Sul e -0,08 para o Norte ). encontrada na região Sul ( 45,7 ).
A nota mínima nesta parte foi zero em todas as regiões, sem exceção.
O Gráfico 3.5 propicia a avaliação do desempenho dos estudantes no Componente de Formação Geral a partir do histograma da distribuição das notas correspondentes.
A distribuição é unimodal, com moda em ( 40 ; 50 ], enquanto na prova como um todo a moda foi alcançada no intervalo ( 50 ; 60 ].
Nota -se, ainda, que no Gráfico 3.5 as notas apresentam uma maior dispersão do que no Gráfico 3.1 ( distribuição das notas da prova ), confirmado pela comparação dos desvios padrões : 13,9 para a nota da prova como um todo e 14,7 para o Componente de Formação Geral.
Para o Componente de Formação Geral, o coeficiente de assimetria da distribuição das notas dos estudantes é negativo ( -0,16 ).
Por ser um coeficiente pequeno, a distribuição apresenta uma leve concentração à direita e cauda maior à esquerda, porém podemos considerar que a distribuição é aproximadamente simétrica.
Em todas as Grandes Regiões os histogramas também possuem assimetria levemente negativa ( entre -0,20 para o Sul e -0,08 para o Norte ).
Nos Gráficos 3.6, 3.7 e 3.8 são apresentadas as informações referentes ao desempenho dos concluintes no Componente de Formação Geral, em diferentes agregações : Grande Região do país, Categoria Administrativa e Organização Acadêmica.
Observa -se pelo Gráfico 3.6 que existe diferença estatisticamente significativa ao nível de 95 \% entre a maior média das notas no Componente de Formação Geral, obtida na região Sul ( 45,3 ) e a menor, obtida na região Norte ( 42,8 ).
Vemos que o intervalo de confiança mais largo é o da região Norte ; já o intervalo mais estreito é observado tanto na região Nordeste quanto na região Sudeste.
Este fato está relacionado, também, com o tamanho da população envolvida, menor na região Norte do que na Nordeste e na Sudeste.
No Gráfico 3.7, que representa as notas médias no Componente de Formação Geral segundo Categoria Administrativa do país, observa -se que existe diferença estatisticamente significativa entre as médias.
Assim como ocorreu para as notas da prova como um todo, para o Componente de Formação Geral, os concluintes de Enfermagem, das IES Públicas ( 49,0 ) obtiveram uma média maior do que os das IES Privadas ( 43,6 ).
Considerando -se o tipo de Organização Acadêmica, nota -se, no Gráfico 3.8, que existe uma diferença estatisticamente significativa entre a maior e a menor média.
Nas Universidades, a média ( 46,4 ) é mais elevada do que em Centros Universitários ( 43,4 ) e Faculdades ( 43,3 ). 3.1.3 Estatísticas Básicas do Componente de Conhecimento Específico A Tabela 3.3 apresenta as estatísticas básicas referentes ao Componente de Conhecimento Específico da área de Enfermagem.
A média do desempenho dos alunos do Brasil como um todo foi 53,4.
A maior média foi obtida na região Sul ( 58,0 ), e a menor, na região Norte ( 50,7 ).
As demais médias foram : 52,9 na região Nordeste, 53,4 na região Sudeste e 52,1 na região Centro-Oeste.
Quanto à variabilidade das notas, o desvio padrão do Brasil como um todo foi 15,9, sendo o maior desvio padrão observado na região Nordeste ( 16,3 ) e o menor na região Sul ( 14,3 ).
O desvio padrão das demais regiões foi : 14,8 na região Norte, 15,9 na região Sudeste e 15,6 na região Centro-Oeste.
A mediana das notas dos estudantes de todo o Brasil foi 54,3.
A maior mediana ocorreu na região Sul ( 59,4 ) e a menor na região Norte ( 51,5 ).
As demais medianas foram : 53,8 na região Nordeste, 54,3 na região Sudeste e 52,6 na região Centro-Oeste.
A nota máxima do Brasil como um todo foi 95,5, sendo obtida por pelo menos um aluno da região Sudeste.
As demais notas máximas foram : 91,2 na região Norte, 95,0 na região Nordeste, 94,8 na região Sul e 93,0 na região Centro-Oeste.
A nota mínima foi zero em todas as regiões, sem exceção.
Assim como os Gráficos 3.1 e 3.5, o Gráfico 3.9, apresentado a seguir, proporciona uma avaliação do desempenho de concluintes em relação ao Componente de Conhecimento Específico com um histograma da distribuição das notas correspondentes.
Esta também é uma distribuição unimodal, e o grupo modal é o ( 50 ; 60 ], acima do grupo modal para a Formação Geral.
O coeficiente de assimetria da distribuição das notas do Componente de Conhecimento Específico é negativo ( -0,29 ).
Nota -se pelo histograma ( Gráfico 3.9 ) que esta distribuição apresenta uma leve concentração à direita e cauda maior à esquerda.
Em todas as Grandes Regiões os histogramas também possuem assimetria negativa ( entre -0,51 para o Sul e -0,13 para o Norte ).
Os Gráficos 3.10, 3.11 e 3.12 apresentam, respectivamente, uma comparação dos resultados em relação à Grande Região do país, à Categoria Administrativa e à Organização Acadêmica, agora levando em conta o desempenho de alunos no Componente de Conhecimento Específico da prova.
Pelo Gráfico 3.10, observa -se que existe diferença estatisticamente significativa ao nível de 95 \% entre a maior média, obtida na região Sul ( 58,0 ) e a menor, obtida na região Norte ( 50,7 ).
Quanto à Categoria Administrativa ( Gráfico 3.11 ), observa -se um comportamento semelhante àquele da Parte de Formação Geral e à prova como um todo, ou seja, existe diferença estatisticamente significativa entre as médias das IES Públicas e Privadas, sendo que a maior média foi obtida por alunos de IES Públicas de ensino ( 62,5 ).
Quanto ao Gráfico 3.12, observa -se, mais uma vez, que existe diferença estatisticamente significativa ao nível de 95 \% entre as notas no Componente de Conhecimento Específico de Universidades e dos demais tipos de Organização Acadêmica.
Sendo que a média dos concluintes das Universidades ( 57,5 ) foi maior do que de Centros Universitários ( 51,9 ) e de Faculdades ( 50,3 ). 3.2 ANÁLISE DAS QUESTÕES OBJETIVAS 3.2.1 Componente de Formação Geral A Tabela 3.4 apresenta as estatísticas básicas relativas às oito questões objetivas do componente da prova que abrange a Formação Geral dos estudantes.
A média do Brasil foi 46,8.
A menor média foi encontrada na região Norte ( 45,1 ) e a maior na região Sul ( 48,0 ).
As demais médias foram 46,9 na região Nordeste, 47,1 na região Sudeste e 46,0 na região Centro-Oeste.
O desvio padrão do Brasil foi 17,4, sendo o maior desvio padrão encontrado na região Centro-Oeste ( 17,7 ) e o menor na região Sul ( 16,8 ).
Os demais desvios foram : 17,3 na região Norte, 17,5 na região Nordeste, e 17,2 na região Sudeste.
As mesmas medianas ( 50,0 ), notas máximas ( 100,0 ) e notas mínimas ( 0,0 ), nas questões objetivas do Componente de Formação Geral, foram obtidas em todas as regiões, sem exceção.
A Tabela 3.5 apresenta o índice de facilidade e o índice de discriminação ( ponto bisserial ) para cada uma das questões objetivas do Componente de Formação Geral.
Quanto ao índice de facilidade, foram usadas as seguintes cores para diferenciar o nível de dificuldade da questão :  Azul para as questões classificadas com índice muito fácil ( índice \&gt ; =0,86 ), verde para as questões classificadas com índice fácil ( 0,61 a 0,85 ), amarelo para as questões classificadas com médio ( 0,41 a 0,60 ), vermelho para as questões classificadas com difícil ( 0,16 a 0,40 ) e roxo para as questões classificadas com muito difícil ( \&lt ; =0,15 ).
Já quanto ao índice de discriminação, foram usadas as seguintes cores para qualificar a questão :  As questões classificadas com índice fraco receberam a cor vermelha ( índice \&lt ; =0,19 ), as classificadas com médio receberam a cor amarelo ( 0,20 a 0,29 ), as classificadas com bom receberam a cor verde ( 0,30 a 0,39 ) e as classificadas com muito bom ( \&gt ; =0,40 ) receberam a cor azul.
As questões objetivas do Componente de Formação Geral, segundo o índice de facilidade, foram assim avaliadas : das oito questões, uma teve o índice de facilidade classificado como muito fácil e duas como fácil.
Uma questão foi tida com grau médio de dificuldade, por ter índice de acertos situado na faixa entre 0,41 e 0,60 ( de 41,0 \% a 60,0 \% de acertos ).
Quatro questões foram consideradas de dificuldade difícil, situando -se no intervalo entre 0,16 e 0,40 do índice de facilidade, ou seja, houve entre 16,0 \% e 40,0 \% de acertos.
Nenhuma questão foi considerada como muito difícil segundo o índice de facilidade.
Como já comentado, para análise das questões objetivas relativas à Formação Geral segundo o poder de discriminação, utilizou -se, o índice de discriminação ponto bisserial.
A Tabela 3.5 apresenta o índice de facilidade e o índice de discriminação ( ponto bisserial ) para cada uma das questões objetivas do Componente de Formação Geral.
Quanto ao índice de facilidade, foram usadas as seguintes cores para diferenciar o nível de dificuldade da questão :  Azul para as questões classificadas com índice muito fácil ( índice \&gt ; =0,86 ), verde para as questões classificadas com índice fácil ( 0,61 a 0,85 ), amarelo para as questões classificadas com médio ( 0,41 a 0,60 ), vermelho para as questões classificadas com difícil ( 0,16 a 0,40 ) e roxo para as questões classificadas com muito difícil ( \&lt ; =0,15 ).
Já quanto ao índice de discriminação, foram usadas as seguintes cores para qualificar a questão :  As questões classificadas com índice fraco receberam a cor vermelha ( índice \&lt ; =0,19 ), as classificadas com médio receberam a cor amarelo ( 0,20 a 0,29 ), as classificadas com bom receberam a cor verde ( 0,30 a 0,39 ) e as classificadas com muito bom ( \&gt ; =0,40 ) receberam a cor azul.
As questões objetivas do Componente de Formação Geral, segundo o índice de facilidade, foram assim avaliadas : das oito questões, uma teve o índice de facilidade classificado como muito fácil e duas como fácil.
Uma questão foi tida com grau médio de dificuldade, por ter índice de acertos situado na faixa entre 0,41 e 0,60 ( de 41,0 \% a 60,0 \% de acertos ).
Quatro questões foram consideradas de dificuldade difícil, situando -se no intervalo entre 0,16 e 0,40 do índice de facilidade, ou seja, houve entre 16,0 \% e 40,0 \% de acertos.
Nenhuma questão foi considerada como muito difícil segundo o índice de facilidade.
Como já comentado, para análise das questões objetivas relativas à Formação Geral segundo o poder de discriminação, utilizou -se, o índice de discriminação ponto bisserial.
Nesta análise, as questões foram assim avaliadas : cinco das oito questões apresentaram índices acima de 0,40 e, assim, foram classificadas com índice muito bom para esse grupo de alunos ; duas questões tiveram índice de discriminação bom, entre 0,30 e 0,39 ; e uma questão teve índice de discriminação médio, entre 0,20 e 0,29.
Nenhuma questão foi classificada com um índice fraco e, portanto, nenhuma questão foi eliminada pelo critério ponto bisserial.
O índice de facilidade variou de 0,16 a 0,94, e o de discriminação, de 0,29 a 0,47.
As questões com índice de discriminação muito bom, de números 3, 4, 6, 7 e 8, ficaram distribuídas entre quase todos os diferentes níveis de facilidade desse conjunto : duas ( questões 4 e 8 ) classificadas na categoria fácil, uma na categoria médio ( questão 6 ) e duas ( questões 3 e 7 ) classificadas na categoria difícil do índice de facilidade.
Em particular, a questão 3 apresentou maior poder discriminatório, com índice 0,47, e foi a quinta mais fácil, com uma proporção de 0,38 de acertos.
A questão de número 5 apresentou índice de facilidade 0,16, ou seja, um quantitativo de 16,0 \% de estudantes conseguiu resolvê -la, dentro do universo de participantes.
Além disso, seu índice de discriminação foi bom, com valor igual a 0,30.
O Gráfico 3.13, para exemplificar, analisa o comportamento da questão de número 3 de Formação Geral.
Trata -se da quinta questão mais fácil e a que obteve o maior índice de discriminação dessa parte da prova.
Neste gráfico, cada uma das cinco curvas representa o percentual de respostas em determinada alternativa da questão, em função da nota dos estudantes nesta parte da prova ( Formação Geral/Múltipla Escolha ), antes de possíveis eliminações pelo critério do ponto bisserial.
A curva em vermelho corresponde à alternativa E, a correta para esta questão.
Assim, observa -se que entre os estudantes com notas mais baixas, nessa parte do exame, a situação mais frequente foi a escolha de uma das alternativas incorretas : a alternativa A ( em azul ) ou C ( em preto ).
Na medida em que a nota aumenta, indicando desempenho melhor nesta parte da prova, aumenta concomitantemente a proporção de estudantes que selecionaram a alternativa correta E, atingindo 100 \% para os alunos com 8 acertos.
Essa análise permite verificar como a questão discriminou os grupos de desempenho, justificando o alto índice obtido na questão.
Os gráficos relativos às demais questões de Formação Geral constam do Anexo I. 3.2.2 Componente de Conhecimento Específico A Tabela 3.6 apresenta as estatísticas básicas em relação às questões objetivas do Componente de Conhecimento Específico da prova, por Grande Região.
A média do Brasil deste componente foi de 57,6.
A menor média foi observada na região Norte ( 54,6 ) e a maior na região Sul ( 62,6 ).
O desvio padrão de todo o Brasil foi 16,9, sendo o menor desvio padrão encontrado na região Sul ( 15,4 ) e o maior na região Nordeste ( 17,3 ).
A mediana de todo o Brasil foi 57,1, assim como para todas as regiões do Brasil, exceto para a região Sul ( 61,9 ).
A nota máxima da prova ( 100,0 ) foi obtida, nas questões objetivas do Componente de Conhecimento Específico, por pelo menos um aluno em todas as regiões do Brasil, exceto na Norte ( 95,2 ).
Em todas as regiões, sem exceção, a nota mínima foi zero.
A Tabela 3.7 apresenta o índice de facilidade e o índice de discriminação ( ponto bisserial ) para cada uma das questões objetivas do Componente de Conhecimento Específico.
Para facilitar a diferenciação das questões usaram -se as mesmas cores da Tabela 3.5 para as diferentes classificações dos índices de facilidade e de discriminação.
Dentre as questões objetivas da parte da prova relativa ao Componente de Conhecimento Específico, apenas a questão 31 foi anulada pela Comissão.
Desse modo, a classificação quanto ao índice de facilidade foi estabelecida com base em 26 questões restantes.
A partir dos índices obtidos, pode -se concluir que um pouco mais de 1/3 das 26 questões objetivas da prova foi considerada como fácil, dez ao todo.
As 16 questões restantes foram classificadas com índice de facilidade da seguinte maneira : uma questão foi muito fácil, nove foram médias e seis foram difíceis.
Já quanto aos índices de discriminação das questões objetivas do Componente de Conhecimento Específico da prova, tem -se como resultado a seguinte classificação : 11 das 26 questões foram consideradas como boas, enquanto seis delas tiveram índice de discriminação muito bom.
Assim, para 17 das 26 questões, os índices de discriminação foram bons ou muito bons.
Dentre as demais, quatro delas foram classificadas como médias e outras cinco como fracas, sendo nove, por conseguinte, a quantidade de questões nos dois patamares mais baixos de discriminação.
Constata -se, assim, que a prova – no que se refere ao Componente de Conhecimento Específico – possuía capacidade mediana de discriminar entre aqueles que dominam ou não o conteúdo.
Dentre as questões que alcançaram os maiores índices de discriminação, as de números 9, 10, 13, 15, 24 e 28 classificadas com índice muito bom, igual ou superior a 0,40, quatro delas ( questões 10, 13, 24 e 28 ) foram classificadas na categoria fácil e as outras duas ( questões 9 e 15 ) como médio, quanto ao índice de facilidade.
A questão de número 33 foi a mais difícil dentre as 26 questões específicas consideradas, com baixo índice de facilidade, apenas 17,0 \% de acertos.
Essa questão apresentou poder discriminatório igualmente baixo, 0,08, o que comprova ter sido esta a mais difícil para os estudantes.
Destacam -se, também, as questões 30, 22 e 25.
A questão 30, por ter índice de facilidade 0,26, o que, em termos percentuais, corresponde a 26,0 \% de estudantes que responderam acertadamente, obtendo, ainda, 0,10 de índice de discriminação.
Já a questão 22, com índice de facilidade 0,27, o que, em termos percentuais, corresponde a 27,0 \% de estudantes que responderam acertadamente, obteve 0,15 de índice de discriminação.
Por fim, a questão 25, com índice de facilidade 0,34, o que, em termos percentuais, corresponde a 34,0 \% de estudantes que responderam acertadamente, obteve 0,12 de índice de discriminação.
Tais questões foram, portanto, quatro das mais difíceis da prova.
As questões 22, 25, 30 e 33 foram eliminadas do cômputo da nota final pelo critério do ponto bisserial.
Além destas quatro, a questão 27, com índice fraco de discriminação, também não foram computadas.
A título de exemplo das análises do comportamento das questões objetivas, o Gráfico 3.14 analisa a questão 28 do Componente de Conhecimento Específico.
Esta foi uma das questões classificadas como fácil nesta prova, apresentando índice de facilidade 0,61, ou seja, 61,0 \% dos estudantes assinalaram acertadamente a opção C, correspondente ao gabarito.
Seu índice de discriminação foi igual a 0,43, classificado como muito bom.
Neste gráfico, cada uma das cinco curvas representa o percentual de respostas em determinada alternativa da questão 28, em função da nota dos estudantes nesta parte da prova, antes de possíveis eliminações de questões pelo critério do ponto bisserial.
A alternativa correta C, representada no gráfico pela curva em preto, foi escolhida em maiores proporções pelos alunos com desempenho melhor nesta parte da prova.
Já as alternativas incorretas, também denominadas distratores, foram selecionadas principalmente por aqueles com notas mais baixas.
A proporção de alunos que selecionou a resposta correta C aumenta gradativamente, chegando a atingir 100 \% para aqueles que acertaram pelo menos 24 questões, enquanto a proporção dos que escolheram alternativas incorretas decai, a partir da primeira nota não nula, como função da nota nesta parte da prova.
Os gráficos relativos às demais questões do Componente de Conhecimento Específico constam do Anexo I. 3.3 ANÁLISE DAS QUESTÕES DISCURSIVAS 3.3.1 Componente de Formação Geral As análises dos resultados de desempenho dos estudantes de Enfermagem nas duas questões discursivas relativas à Formação Geral encontram -se na Tabela 3.8 e no Gráfico 3.15.
Na Tabela 3.8 observa -se que as notas médias foram mais baixas nesse conjunto de questões do que no das objetivas.
Os estudantes de todo o Brasil obtiveram, em Formação Geral, média 46,8 nas questões objetivas e 41,2 nas questões discursivas.
A mediana também confirma o pior desempenho dos alunos de todo o Brasil nas questões discursivas do Componente de Formação Geral.
Enquanto essa estatística foi de 46,5 para questões discursivas, para as questões objetivas a mediana foi de 50,0.
Pode -se, também, notar um aumento do desvio padrão de 17,4 nas questões objetivas do Componente de Formação Geral dos alunos de todo o Brasil, para 22,0 nas questões discursivas do mesmo componente.
Como já comentado, a mediana de todo o Brasil, neste componente, foi 46,5, sendo esta encontrada na região Sudeste, mais alta para as regiões Nordeste e Sul ( 47,0 ) e mais baixa para as regiões Norte ( 42,0 ) e Centro-Oeste ( 46,0 ).
A maior nota máxima foi obtida, nas questões discursivas do Componente de Formação Geral, por pelo menos um aluno da região Nordeste ( 97,0 ), enquanto a menor nota máxima foi obtida nas regiões Norte e Sul ( 87,5 ).
Em todas as regiões a nota mínima foi zero, sem exceção.
O Gráfico 3.15 representa a distribuição das notas nas questões discursivas no Componente de Formação Geral.
A moda desta distribuição ocorre no intervalo ( 50 ; 60 ].
Devido à grande quantidade de notas zero e à alta frequência de alunos que deixaram este tipo de questão em branco o intervalo [ 0 ; 10 ] se apresenta como uma segunda moda.
A distribuição possui assimetria à esquerda, coeficiente de assimetria -0,61.
Em todas as regiões o coeficiente de assimetria também é negativo, o maior em módulo, -0,72, na região Sul, e o menor na região Norte ( -0,42 ).
A avaliação das questões discursivas de Formação Geral considerou em separado o conteúdo ( \textbf{peso} 0,8 na nota ) e o desempenho em língua portuguesa ( \textbf{peso} 0,2 ).
Na sequência, os resultados verificados para cada uma das questões discursivas de Formação Geral com relação ao conteúdo serão apresentados, estabelecendo -se relações com os temas abordados em cada uma delas.
Os comentários da Banca de docentes corretores a respeito do observado na correção das respostas dos estudantes, suas impressões e conclusões serão apresentados junto à análise de cada questão.
Em seguida será feita uma análise do desempenho de Língua Portuguesa.
Os comentários da Banca de docentes corretores serão apresentados para o conjunto de questões.
Cumpre esclarecer que, tendo em vista que as questões discursivas de Formação Geral são padronizadas, ou seja, constam de todas as provas, os comentários da Banca são os mesmos para todas as carreiras acadêmicas, sendo direcionados a todos os estudantes que participaram do ENADE/2013.
A seguir, serão analisados os desempenhos da Área de Enfermagem nas duas questões discursivas de Formação Geral do ENADE/2013, comparando os resultados obtidos com comentários para cada questão. 3.3.1.1 Análise de Conteúdo Questão Discursiva 1 do Componente de Formação Geral Os dados de Enfermagem, obtidos a partir das respostas à questão 1, encontram -se na Tabela 3.9 e no Gráfico 3.16.
Nessa questão – de melhor desempenho dentre as duas de Formação Geral – os alunos de todo Brasil tiveram média, 43,6.
A maior média para a questão 1 foi obtida nas regiões Norte e Centro-Oeste ( 45,4 ), e a menor, na região Sudeste ( 41,5 ).
Quanto à variabilidade das notas, o desvio padrão de todo o Brasil foi 26,9.
O menor desvio padrão foi obtido na região Sul ( 26,1 ) e o maior desvio padrão foi obtido na região Sudeste ( 27,1 ).
As mesmas medianas ( 50,0 ), notas máximas ( 100,0 ) e notas mínimas ( 0,0 ), na questão discursiva 1 do Componente de Formação Geral, foram obtidas em todas as regiões, sem exceção.
O Gráfico 3.16 mostra a distribuição das notas na questão discursiva 1 do Componente de Formação Geral.
Observa -se que a frequência modal corresponde aos alunos que deixaram a questão em branco.
Considerando apenas os estudantes que tentaram resolver a questão, a moda ocorre no intervalo ( 40 ; 50 ], com uma moda secundária no intervalo ( 60 ; 70 ].
A distribuição das notas possui assimetria negativa ( -0,42 ), assim como em todas as regiões do Brasil ( entre -0,48 para a Centro-Oeste e -0,36 para as regiões Norte e Sudeste ). 3.3.1.2 Comentários sobre a correção de Conteúdo das respostas à Questão Discursiva 1 O enunciado restou claro, elaborado com linguagem técnica, estando inserido nas Diretrizes gerais, com nível médio de profundidade, já que abordava tema vinculado à Área de formação da maior parte dos examinandos.
A divisão das perguntas também ficou clara indicando, passo a passo, o caminho de resposta a seguir, com base no texto.
Verificou -se dificuldade quanto à análise dos programas de governo existentes que contribuem para melhoria do saneamento básico.
Nesse ponto, os alunos referiram, com muita frequência, a ações relacionadas à Saúde, tais como : médico da família, ações educativas, ações ambientais, melhoria nas condições de habitação ( sem necessariamente vincular a programas de governo federal ou estaduais ou municipais ).
Os programas citados foram : Programa Saúde da Família ( PSF ), Programa de \textbf{Aceleração} do Crescimento ( PAC ) e Minha Casa Minha Vida e, ainda, o SUS.
Apesar de uma minoria de examinandos ter indicado o PAC do Governo Federal, a maior parte deles apresentou alguma solução para a questão do saneamento básico.
Os estudantes indicaram, com facilidade, a questão da necessidade do tratamento da água como item essencial do saneamento básico.
Muitos concluintes, através da sua experiência local, citaram situações em que as crianças caminhavam de pés no \textbf{chão} perto de rotas de esgoto a céu aberto.
Também criticaram a inércia governamental em não combater as invasões ilícitas, em locais inadequados, sem estrutura de água potável e também de saneamento básico.
Diante das \textbf{observações} preliminares pode ser extraída uma classificação das respostas obtidas que ficou estabelecida em fracas, médias e boas.
As notas consideradas fracas ( notas de 0 a 3,0 ) foram conferidas àqueles que não conseguiram relacionar o saneamento básico com a importância do tratamento de água, os projetos para resolver o problema e nem indicaram soluções para a questão.
As respostas que tiveram média avaliação ( notas de 3,5 a 7,0 ) tocaram, de algum modo, nos pontos indicados, mesmo que de forma insuficiente.
Já as respostas consideradas boas ( notas de 7,5 a 10,0 ) preencheram os requisitos do padrão de resposta, não logrando nota máxima, na maior parte dos casos, por ausência de indicação de programa governamental, mas sugerindo soluções técnicas para o problema.
As respostas mais comuns foram as consideradas medianas e que indicaram, com frequência, que o problema do saneamento básico estaria nas classes de renda mais desfavorecida e nas regiões consideradas como periferias das grandes cidades.
Também foram frequentes respostas que citaram a necessidade de tratamento de água, acompanhada da necessidade de maiores investimentos nessa área. 3.3.1.3 Análise de Conteúdo da Questão Discursiva 2 do Componente de Formação Geral A Tabela 3.10 mostra que o desempenho dos estudantes na questão 2 ( média 35,8 ) foi inferior ao obtido na questão de número 1 ( média 43,6 ).
A região Sul foi aquela onde a média, nessa questão, foi maior ( 37,6 ), e a de menor média foi a região Norte ( 29,7 ).
Quanto à variabilidade das notas, o desvio padrão de todo o Brasil foi 26,9, o mesmo obtido na O maior desvio nessa questão foi obtido na região Nordeste ( 27,9 ), enquanto o menor foi obtido na região Norte ( 24,6 ).
A mediana foi 40,0 em todas as regiões, exceto na região Norte ( 30,0 ).
As notas máximas ( 100,0 ) foram as mesmas em todas as regiões do Brasil, assim como as notas mínimas ( 0,0 ), sem exceção.
O Gráfico 3.17 mostra a distribuição das notas na questão discursiva 2 do Componente de Formação Geral.
Neste gráfico, desconsiderando -se a grande quantidade de alunos que deixaram a questão 2 em branco ( em torno de 23 \% ), observa -se a moda no intervalo ( 40 ; 50 ].
O coeficiente de assimetria das notas nesta questão é negativo ( -0,03 ), mas quase nulo, evidenciando o balanço dos alunos que deixaram a questão em branco ( e receberam nota zero ) e os demais.
O mesmo ocorre para todas as regiões da Brasil ( entre -0,17 na região Sul e -0,04 na Centro-Oeste ), exceto para a região Norte ( 0,36 ) que possui coeficiente de assimetria positivo. 3.3.1.4 Comentários sobre a correção de Conteúdo das respostas à Questão Discursiva 2 O tema abordado no enunciado, apesar da complexidade, é extremamente atual.
Essa atualidade facilitou a elaboração de alguma resposta, permitindo que todos os que se dispuseram a responder pudessem apresentar, pelo menos, alguns dos elementos integrantes do padrão de resposta.
A correção, seguindo o padrão de resposta, avaliou a capacidade de elaborar um texto dissertativo que abordasse o tema sob três pontos de vista : o papel da tecnologia digital na atualidade ; a garantia dos direitos do cidadão e do Estado diante do avanço tecnológico ; e o problema da segurança digital para o indivíduo, para o Estado e para as organizações ( empresas e/ou instituições nacionais ou internacionais ).
Dada a amplitude de respostas possíveis, os alunos tiveram desempenho satisfatório.
Identificamos que os concluintes, basicamente, trilharam dois caminhos : a ) a análise da influência das inovações tecnológicas, notadamente através das redes sociais, numa perspectiva essencialmente privada, daí defluindo as questões da privacidade e da segurança ; b ) uma perspectiva de segurança nacional-estatal com criticas à atuação americana de invasão de privacidade dos dados governamentais, e à inercia do governo brasileiro em reagir ao ataque, considerado, pela maioria esmagadora dos concluintes, como indevido e antiético.
Poucos foram os que uniram as duas facetas da questão : o público e o privado.
Na maior parte das provas, o concluinte indicou a ausência de segurança e de privacidade decorrentes das próprias atividades individuais.
Muitos respondentes citaram que, de forma voluntária, os próprios cidadãos enfraquecem a proteção à privacidade, pois divulgam sua vida pessoal, aspectos de seu cotidiano, por meio de imagens e do fornecimento de dados pessoais.
No aspecto segurança, foram inúmeras as indicações de acessos indevidos a dados pessoais com a perpetração de fraudes de todos os tipos, notadamente as bancárias e comerciais, com a aquisição de bens à revelia dos titulares dos dados.
Também no tema segurança, muitos indicaram o aspecto positivo decorrente da vigilância generalizada, não somente através da rede de computadores, mas também pela quantidade de \textbf{câmeras} existentes, que contribuem para a identificação de criminosos procurados ou dos que realizam crimes pela internet.
Na perspectiva pública, a esmagadora maioria indicou repulsa pelo comportamento americano de invasão dos dados do governo brasileiro, reclamando reação do mesmo.
Alguns poucos justificaram a intervenção americana pela “ natural ” \textbf{disputa} de poder pelos Estados imperialistas. 3.3.1.5 Análise de Língua Portuguesa das Questões Discursivas do Componente de Formação Geral Os dados de Enfermagem, obtidos a partir das respostas às questões discursivas do Componente de Formação Geral no que tange à Língua Portuguesa, encontram -se na Tabela 3.11 e no Gráfico 3.18.
Nesse aspecto, os alunos de todo Brasil tiveram média, 43,1.
A maior média com respeito à Língua Portuguesa foi obtida na região Norte ( 46,4 ), e a menor, na região Sudeste ( 41,1 ).
Quanto à variabilidade das notas, o desvio padrão de todo o Brasil foi 22,8.
O menor desvio padrão foi obtido na região Norte ( 21,5 ) e o maior desvio padrão foi obtido na região Sudeste ( 23,2 ).
A mediana do Brasil como um todo da nota de Língua Portuguesa foi 50,0, sendo esta encontrada também na região Nordeste.
A mediana foi mais alta para as regiões Norte e Sul ( 52,5 ) e mais baixa para as regiões Sudeste e Centro-Oeste ( 47,5 ).
A nota máxima foi 90,0 para a maioria das regiões, exceto para as regiões Sudeste ( 87,5 ) e Sul ( 85,0 ).
Além disso, a nota mínima foi zero em todas as regiões do país, sem exceção.
O Gráfico 3.18 mostra a distribuição das notas de Língua Portuguesa do Componente de Formação Geral.
Observa -se que a maior frequência corresponde aos alunos que obtiveram nota no intervalo ( 50 ; 60 ], correspondendo à moda da distribuição.
Cerca de 15 \% dos alunos deixaram a prova em branco, e esta situação corresponde a uma moda secundária.
A distribuição das notas possui assimetria negativa ( -0,76 ), assim como em todas as regiões do Brasil ( entre -0,94 para o Sul e -0,73 para o Nordeste e o Sudeste ). 3.3.1.6 Comentários sobre a correção das respostas de Formação Geral com respeito à Língua Portuguesa Um aspecto interessante a ser destacado é a diferença de desempenho entre a questão 1 e a questão 2, principalmente nos aspectos textuais – os textos produzidos para a questão 1 apresentaram resultados melhores do que os produzidos para a questão 2.
Uma hipótese que provavelmente explique essa discrepância é o fato de o tema “ saneamento básico ” ser mais relevante para profissionais de saúde do que o tema “ segurança e privacidade ”, constituindo tema trabalhado no currículo escolar.
Prova disso é que quase todos os textos da questão 1 citam termos específicos da área de saúde, como “ contaminação feco-oral ”, “ higiene corporal ”, “ ingestão de alimentos ou água contaminada por vetores biológicos ”.
Os resultados da avaliação correspondem aos seguintes aspectos observados em cada competência : a ) aspectos ortográficos : o desempenho dos participantes ficou concentrado, majoritariamente, nos níveis 3 ( até 3 desvios ) e 2 ( de 4 a 7 desvios ), devido ao grande índice de desvios de acentuação, contrabalançada pelo pequeno número de desvios em grafia.
Foi reduzido o número de textos com ausência de desvios ( enquadrados no nível 4 ).
Alguns textos, com pior desempenho, foram enquadrados no nível 1 ( de 8 a 12 desvios ), enquanto o nível zero foi atribuído a pouquíssimos casos, já que ele revela falta absoluta de domínio das convenções ortográficas.
Observou -se, portanto, que existe uma diferença muito grande de desempenho nos dois aspectos analisados : baixíssimo índice de desvios ortográficos e grande índice de desvios de acentuação.
Em vários casos, ocorre ausência completa de acentuação gráfica.
Os resultados revelam, portanto, que a tendência dominante entre os universitários brasileiros é a eliminação da acentuação gráfica, talvez motivada pela ausência de esclarecimento dos meios de comunicação, das autoridades e das escolas sobre as decisões do Acordo Ortográfico da Língua Portuguesa de 1990.
Os casos mais sistemáticos de ausência de acentuação são :  palavras proparoxítonas ( “ politicas ”, “ publicas ”, “ juridicas ” ) ;  palavras paroxítonas terminadas em ditongo crescente ( “ varios ”, “ Brasilia ”, “ familia ” ) ;  palavras oxítonas ( “ ninguem ”, “ esta ”, “ ate ” ).
Quanto ao domínio das convenções relativas à grafia das palavras, são pontuais os problemas encontrados ( “ infecciozas ”, “ obtos ”, “ hipedemia ”, “ saniamento ”, “ cituação ”, “ abitantes ” ).
Os únicos desvios recorrentes foram : a ausência de segmentação nos casos de verbo + pronome oblíquo “ se ”, combinada, eventualmente, à alteração da grafia do pronome ( “ percebese ”, “ percebece ”, “ preferece ” ) ; a grafia da forma verbal de pretérito imperfeito do subjuntivo como se fosse uma ocorrência de verbo + pronome oblíquo ( “ distribui -se ” por “ distribuísse ” ) Vale observar, também, que, ao contrário do que se esperava, não apareceram abreviaturas próprias do “ internetês ”, ou seja, dos hábitos de comunicação escrita adquiridos pelo uso de redes sociais e e-mails. b ) aspectos textuais : o desempenho dos participantes ficou concentrado, majoritariamente, nos níveis 3 e 2, devido à grande ocorrência de problemas de estruturação textual.
Foi muito baixo o número de textos enquadrados no nível 4 por não apresentar qualquer problema estrutural.
Alguns textos com pior desempenho foram enquadrados no nível 1, enquanto o nível zero foi atribuído a pouquíssimos casos, já que ele revela a existência de um texto sem articulação e com comprometimento do sentido.
Observou -se que a grande maioria dos participantes não distribuiu as ideias em parágrafos, talvez devido ao pequeno número de linhas disponibilizadas para a resposta da questão ou, quem sabe, pela suposição de que não seria necessária essa divisão por não se tratar de um texto no modelo de uma redação dissertativo-argumentativa, como solicitado nos vestibulares.
Essa competência é a que se revela como a mais problemática entre os participantes, porque são muitos os problemas observados, desvios acumulados durante toda a formação do estudante e que não se resolvem com um estudo autodidata, como acontece com regras ortográficas ou morfossintáticas : sequência justaposta de ideias sem encaixamentos sintáticos ; redução drástica de estruturas subordinadas, ao lado do aumento na frequência de estruturas coordenadas e absolutas ; redução no uso de conectores para expressar relações lógicas essenciais à construção do texto, substituídas pela exigência de inferência por parte do interlocutor para suprir a sua ausência ; emprego equivocado de operadores que não estabelecem relações lógicas coerentes entre ideias do texto ; emprego inadequado do pronome relativo ( com omissão da preposição ou a utilização de pronome inadequado, como “ onde ” ) ; repetição exaustiva de termos sem a utilização de procedimentos mais sofisticados de substituição ( hiperonímias, hiponímias, nominalizações, expressões metafóricas ) ; frases fragmentadas que comprometem a estrutura lógico-gramatical ; frases formadas apenas por oração subordinada, sem oração principal.
Um importante aspecto a destacar é o baixíssimo desempenho de uma parte dos participantes em relação à estrutura formal do texto produzido, o que é extremamente preocupante ao se levar em conta que são graduandos em fase final de formação.
São frequentes os casos de desvios de estruturação frasal, com uso inadequado ou ausência de conectivos entre parágrafos e entre frases.
Em uma parte dos textos, falta um mínimo de textualidade e de domínio do registro padrão da língua.
Na verdade, observam -se relações linguísticas quase agramaticais, como as estabelecidas pela sequência de gerúndios sem o apoio de um ponto de partida para a organização das informações gramaticais e semânticas.
Quanto à utilização dos mecanismos de referenciação, dois aspectos devem ser destacados : a ocorrência, em uma boa parte dos textos, de repetições de palavras ou expressões sem a utilização de termos sinônimos ou pronomes, como seria adequado ; e a alternância entre os pronomes demonstrativos este/ esse, e suas formas flexionadas, sem a observância de um critério de utilização.
Esse último aspecto tem se generalizado na modalidade oral da língua e, cada vez mais, tem se tornado uma constante na modalidade escrita de jornais e revistas : predominância da forma “ esse ” na modalidade oral e da forma “ este ” na modalidade escrita, provavelmente por um processo de hipercorreção. c ) aspectos morfossintáticos e vocabulares : o desempenho dos participantes ficou concentrado, majoritariamente, nos níveis 3 e 2, devido à grande ocorrência de problemas de regência, concordância e uso adequado dos sinais de pontuação.
O nível 4 foi atribuído a um número muito reduzido de textos.
Alguns textos com pior desempenho foram enquadrados no nível 1, enquanto o nível zero foi atribuído a pouquíssimos casos, já que ele revela a existência de um texto sem articulação e com comprometimento do sentido.
Os resultados são muito transparentes em relação aos aspectos mais problemáticos no desempenho dos participantes.
O desvio mais frequente é a falta do sinal indicativo da crase, o que se configura como um problema de regência, verbal ou nominal, conforme o caso – isso revela que o usuário não tem consciência de que, sob a forma do termo “ a ” existe a presença de uma preposição “ a ”, exigida pela regência do termo anterior.
Embora em outros exames, como o Enem, a falta de crase seja \textbf{apenada} em acentuação, nesta correção amostral esse desvio foi considerado no âmbito dos aspectos morfossintáticos.
Outro problema relacionado à regência verbal e à nominal, encontrado frequentemente nas questões, foi a ausência de preposição antes de pronome relativo, processo generalizado na modalidade oral da língua, em situações de registro informal.
Apesar da possibilidade de que essa alteração de regência se generalize no padrão escrito da Língua Portuguesa, como já está ocorrendo até em textos jornalísticos, o não emprego da preposição foi \textbf{apenado} neste processo de avaliação.
Outro desvio muito frequente diz respeito aos processos de concordância verbal e de concordância nominal.
Quanto à concordância de número, observou -se ausência de marca ( com sujeito anteposto ou posposto ) ou uso indevido ( uso inadequado da marca de plural comandado pelo núcleo plural da locução adjetiva, apesar de o substantivo que funciona como núcleo do sintagma nominal estar no singular ).
Uma ocorrência que se destacou foi a ausência de acento circunflexo na forma plural do presente do indicativo do verbo “ ter ”, que foi considerada como um desvio na concordância verbal e não na acentuação gráfica.
Quanto à concordância de gênero, vários casos foram observados, normalmente no âmbito de sintagmas nominais longos, em que o adjetivo está afastado do substantivo.
Quanto à questão da colocação pronominal, foram poucos os casos observados.
Apesar de serem aspectos relacionados à oralidade, concluiu -se que, no registro escrito formal, a maioria dos participantes já incorporou regras como a não introdução da frase por um pronome oblíquo e a próclise na presença de um termo atrator.
Não se adotou, entretanto, o padrão excessivamente formal descrito pelas gramáticas normativas em relação à posição do pronome oblíquo em locuções verbais, já que esse uso está muito distante da prática cotidiana, até em textos mais formais.
Quanto aos aspectos vocabulares, alguns tipos de inadequação foram observados : expressões da oralidade apareceram em algumas respostas, mas sem maior relevância do ponto de vista quantitativo ; seleção vocabular incompatível com o contexto, gerando situações de falta de inteligibilidade ; falta de domínio de vocabulário mais abstrato e de maior complexidade, essencial ao desenvolvimento do texto de base dissertativa.
Quanto à utilização dos sinais de pontuação, observou -se uma grande precariedade nos textos analisados.
Não foi \textbf{apenada} a ausência de vírgula para destacar locuções ou adjuntos adverbiais de pequena extensão deslocados de posição na frase, por ser um uso opcional.
São os seguintes os tipos de problemas encontrados : a ) vírgula : utilização de vírgula para separar o sujeito e o predicado ; ocorrência de apenas uma das vírgulas para separar uma palavra, uma expressão ou uma oração encaixada ; uso de vírgula no lugar do ponto para separar ideias que constituem períodos distintos ; ausência de vírgula para separar enumerações ; ausência de vírgula para separar oração adjetiva explicativa ou utilização inadequada para separar oração adjetiva restritiva ; b ) ponto e vírgula : utilização do ponto e vírgula no lugar de vírgula ; c ) ponto final : ausência de ponto final para separar períodos.
Várias marcas de oralidade foram identificadas, embora não em alta frequência : o uso do pronome relativo “ onde ” como relativo universal, falta de artigo definido antes de substantivo, repetição de palavras por falta de vocabulário, reduções como “ tá ”, “ pra ”, “ pro ”, “ prum ”, expressões informais como “ estão nem aí ”, “ tem que ”, “ levando com a barriga ”, “ não tá nem aí ”, “ a coisa fica séria ”, “ e aí haja dor de cabeça ”. 3.3.2 Componente de Conhecimento Específico Na parte da prova relativa às questões discursivas no Componente de Conhecimento Específico ( Tabela 3.12 ), observa -se que a média foi mais baixa do que nas questões discursivas do Componente de Formação Geral.
Enquanto no Componente de Formação Geral a média para estudantes de Enfermagem de todo o Brasil foi 41,2, na Parte de Conhecimento Específico a média foi 29,8.
A maior média deste componente foi obtida pelos estudantes da região Sul ( 31,7 ), e a menor, pelos da região Nordeste ( 28,5 ).
Quanto à variabilidade das notas, o desvio padrão de todo o Brasil foi 20,2.
O maior desvio padrão foi encontrado na região Nordeste ( 20,5 ), e o menor, nas regiões Norte e Centro-Oeste ( 19,4 ).
O Gráfico 3.19 representa a distribuição das notas nas questões discursivas no Componente de Conhecimento Específico.
Essa distribuição tem moda no intervalo de [ 0 ; 10 ].
Uma moda secundária pode ser reconhecida no intervalo de ( 30 ; 40 ].
O coeficiente de assimetria é positivo para todo o Brasil ( 0,04 ) e para a maioria das regiões ( entre 0,06 para a região Centro-Oeste e 0,15 para a região Norte ), indicando uma assimetria à direita.
Já para as regiões Sudeste ( -0,05 ) e Sul ( -0,12 ) o coeficiente de assimetria é negativo, indicando uma assimetria à esquerda.
A maior nota máxima, nas questões discursivas do Componente de Conhecimento Específico, foi obtida na região Nordeste ( 86,7 ), enquanto a menor nota máxima foi encontrada nas regiões Sul e Centro-Oeste ( 80,0 ).
Além disso, a nota mínima ( 0,0 ) foi obtida por alunos de todas as regiões do Brasil sem exceção.
A mediana do Brasil como um todo foi 30,0, também encontrada na região Centro-Oeste.
As demais medianas foram : 28,3 nas regiões Norte e Nordeste, e 33,3 nas regiões Sudeste e Sul.
O Gráfico 3.19 representa a distribuição das notas nas questões discursivas no Componente de Conhecimento Específico.
Essa distribuição tem moda no intervalo de [ 0 ; 10 ].
Uma moda secundária pode ser reconhecida no intervalo de ( 30 ; 40 ].
O coeficiente de assimetria é positivo para todo o Brasil ( 0,04 ) e para a maioria das regiões ( entre 0,06 para a região Centro-Oeste e 0,15 para a região Norte ), indicando uma assimetria à direita.
Já para as regiões Sudeste ( -0,05 ) e Sul ( -0,12 ) o coeficiente de assimetria é negativo, indicando uma assimetria à esquerda. 3.3.2.1 Análise da Questão Discursiva 3 do Componente de Conhecimento Específico Na questão 3, cujos resultados aferidos encontram -se descritos na Tabela 3.13, a média dos estudantes de todo o Brasil foi 40,5.
A menor média nessa questão foi obtida pelos alunos da região Nordeste ( 36,9 ), enquanto a maior média foi obtida na região Sul ( 45,4 ).
Quanto à variabilidade das notas, o desvio padrão de todo o Brasil foi 30,1.
O maior desvio padrão foi obtido na região Sudeste ( 30,4 ), enquanto o menor foi obtido na região Norte ( 29,2 ).
A nota máxima, 100,0 pontos, foi alcançada por pelo menos um aluno em todas as regiões, sem exceção.
A mediana do Brasil como um todo foi 45,0, assim como nas regiões Norte, Sudeste e Centro-Oeste, mais alta na região Sul ( 55,0 ) e mais baixa na região Nordeste ( 30,0 ).
A nota mínima foi zero para todas as regiões do Brasil, sem exceção.
O Gráfico 3.20 mostra a distribuição das notas na questão discursiva 3 do Componente de Conhecimento Específico.
Essa distribuição apresenta várias modas, mas a principal corresponde às questões em branco.
O coeficiente de assimetria é positivo para todo o Brasil ( 0,02 ) e para a maioria das regiões ( 0,04 nas regiões Norte e Centro-Oeste, e 0,19 na região Nordeste ), indicando uma assimetria à direita.
Já nas regiões Sudeste ( -0,11 ) e Sul ( -0,20 ) o coeficiente de assimetria é negativo, indicando uma assimetria à esquerda. 3.3.2.2 Comentários sobre as respostas à Questão Discursiva 3 O enunciado da Questão 3 era claro e adequado ao conteúdo das Diretrizes Curriculares do Curso de Graduação em Enfermagem que aponta, no perfil do egresso, a capacidade de conhecer e intervir sobre os problemas/situações de saúde-doença mais prevalentes no perfil epidemiológico nacional.
Além disso, aborda um conteúdo relevante, pois está em consonância com a Política Nacional para a Prevenção e Controle do Câncer, publicado em 2013, que considera o perfil epidemiológico vigente no Brasil.
A questão foi considerada de nível médio de profundidade, pela expressão do caso clínico.
O padrão de respostas contemplava parcialmente as respostas apresentadas pelos estudantes.
Diferentes diagnósticos relacionados ao caso apresentado, não contemplados no padrão de resposta, foram descritos pelos estudantes, como por exemplo : conforto prejudicado ; déficit no autocuidado ; integridade tissular prejudicada ; além dos diferentes riscos e vulnerabilidades.
Essa diversidade evidencia certa ampliação do raciocínio crítico, de forma positiva.
A correção também precisou ser flexibilidade quanto à forma redacional dos diagnósticos de enfermagem, pois no Brasil ainda não há uma padronização da taxonomia/linguagem diagnóstica.
No que tange aos critérios de pontuação, os pontos foram distribuídos permitindo a valorização das respostas que evidenciaram maior grau de conhecimento e raciocínio clínico.
A quantidade de respostas com intervenções não relacionadas ao diagnóstico/resultado foi maior do que as convergentes com o diagnóstico ou caso clínico.
Os estudantes, de uma maneira geral, apontaram na assistência de enfermagem atributos como amor, atenção, respeito, paciência, amizade, carinho.
Os diagnósticos de enfermagem mais listados pelos estudantes foram mobilidade física prejudicada/deambulação prejudicada seguido do diagnóstico dor crônica/aguda.
Risco de queda e medo/ansiedade também foram bastante citados.
Contudo, o diagnóstico de enfermagem “ disposição para espiritualidade melhorada ” não foi explicitado na quase totalidade dos documentos corrigidos.
Uma quantidade expressiva de estudantes também citou como diagnóstico de enfermagem “ câncer ”, “ trombose venosa profunda ”, “ fratura ”.
O padrão de respostas não previu os vários riscos que a cliente apresentava e que foram descritos pelos estudantes tais como : risco de integridade da pele prejudicada e risco de comprometimento do padrão circulatório.
Vários estudantes também se referiram à cliente como em pós-operatório de mastectomia.
Muitos, ainda, se posicionaram de maneira enfática quanto ao enfermeiro não poder diagnosticar porque não é profissional médico.
No que diz respeito ao item que solicitava descrever um resultado de enfermagem esperado para o caso clínico, a grande maioria dos estudantes não foi capaz de respondê -lo.
Quanto às intervenções listadas pelos estudantes foi possível verificar que as mesmas não eram específicas para o diagnóstico/problema apresentado.
Intervenções como “ administrar o medicamento prescrito pelo médico ”, “ verificar sinais vitais ”, “ realizar mudança de decúbito ” e “ dar banho no leito ” foram expressivamente encontradas, assim como encaminhamentos não pertinentes.
Em uma apreciação qualitativa é possível afirmar que, em geral, as respostas dos estudantes se relacionavam com as condições clínicas da cliente, contudo, evidenciaram um conhecimento pouco articulado entre teoria e prática, no que se refere a Diagnósticos de Enfermagem.
Foi possível verificar a presença de muitas intervenções de enfermagem contraditórias em uma mesma resposta como, por exemplo, “ auxiliar na deambulação/manter repouso absoluto no leito ” ; “ auxiliar a caminhar até o banheiro/dar banho no leito ; ” “ mudar de decúbito/ retirar a mobília e tapetes do domicílio ”.
Destaca -se, que um número expressivo de estudantes fez referência à necessidade de a cliente ( idosa ) ter acompanhante, citando o Estatuto do Idoso, ou à necessidade da família recorrer aos cuidados de enfermagem com orientações.
Como já comentado, o diagnóstico de enfermagem mobilidade física prejudicada/deambulação prejudicada, seguido do diagnóstico dor crônica/aguda foram os mais citados pelos estudantes.
Risco de queda e medo/ansiedade também está entre os mais citados.
Outras respostas encontradas com muita frequência foram : diagnóstico médico ; ausência de resultado ou resultado esperado inconsistente ; intervenções de enfermagem com pouca especificidade em relação ao caso clínico.
Ao considerarmos as competências e habilidades necessárias ao estudante nessa fase de formação, torna -se importante destacar aspectos observados durante a correção tais como : elevado índice de erros ortográficos e caligrafia ilegível.
A linguagem utilizada por grande parte dos estudantes evidencia um pensamento ainda em construção, com importantes lacunas no raciocínio crítico-reflexivo, tanto na perspectiva da sistematização da assistência de enfermagem, quanto das diferentes taxionomias.
O grande número de notas zero foi atribuído a respostas impossíveis de serem consideradas : diagnóstico médico e não de enfermagem ; ausência de expressão de resultado ou resultado não relacionado com o diagnóstico e, intervenções desencontradas.
Houve também respostas desconsideradas pelo fato de vários estudantes rubricarem, assinarem e/ou informarem o local e a instituição de ensino. 3.3.2.3 Análise da Questão Discursiva 4 do Componente de Conhecimento Específico A Tabela 3.14 contém as informações relativas à questão 4 do conjunto de questões do Componente de Conhecimento Específico.
O desempenho dos estudantes de todo o Brasil nesta questão foi inferior ao desempenho na questão de número 3.
A média geral do Brasil foi 17,0, sendo a menor média registrada na região Sul ( 15,3 ) e a maior na região Sudeste ( 18,3 ).
Quanto à variabilidade das notas, o desvio padrão de todo o Brasil foi 16,8.
O maior desvio padrão foi obtido na região Sudeste ( 17,3 ), enquanto o menor foi obtido na região Sul ( 16,4 ).
A maior nota máxima foi atingida por pelo menos um aluno da região Sul ( 85,0 ), enquanto nas demais regiões a nota máxima foi 75,0 na região Norte, e 80,0 para as regiões Nordeste, Sudeste e Centro-Oeste.
A mediana para o Brasil como um todo foi 15,0, assim como para as regiões Norte e Sudeste, enquanto que para as regiões Nordeste, Sul e Centro-Oeste, a mediana foi 10,0.
A nota mínima foi zero para todas as regiões do Brasil, sem exceção.
O Gráfico 3.21 representa a distribuição de notas na questão discursiva 4, no Componente de Conhecimento Específico.
Como nas outras questões discursivas, o número de estudantes que deixou a questão em branco foi muito grande, correspondendo à moda da distribuição.
Entre aqueles que não deixaram a questão em branco a moda é o intervalo de notas ( 10 ; 20 ].
O coeficiente de assimetria é positivo para todo o Brasil ( 0,78 ) e para todas as regiões ( entre 0,67 para a região Sudeste e 0,96 para a região Sul ), indicando uma assimetria à direita. 3.3.2.4 Comentários sobre as respostas à Questão Discursiva 4 O enunciado da questão apresentava, de maneira clara e objetiva, os sinais e sintomas referentes ao adoecimento por Dengue.
Essa problemática exige um conhecimento específico e aprofundado do enfermeiro para agir com segurança, promovendo um cuidado de enfermagem que, em conformidade com o Ministério da Saúde, 2013, deve ser baseado “ no reconhecimento de elementos clínicos e/ou laboratoriais que podem ser indicativos de gravidade ”.
Os dados da situação em tela permitiram, em parte, que o estudante elaborasse o plano assistencial para a adequada conduta terapêutica de enfermagem requerida para a clínica.
Entretanto, pouquíssimos estudantes desenvolveram o subitem ‘ b ’ da questão em conformidade com essa classificação de risco.
Não houve informações que permitissem avaliar o risco social ( representado por local de residência, tipo de atividade laboral, composição familiar, estado civil ‘ etc. ), mas o enunciado descreveu sinais e sintomas que classificaram a paciente MTS no Grupo C do protocolo do Ministério da Saúde.
Dessa forma, o acolhimento e a triagem pela classificação de risco, como prioridade e responsabilidade do enfermeiro, não foram abordados nos documentos corrigidos, ainda que o texto tenha enunciado, literalmente, o diagnóstico de Dengue Hemorrágica como confirmado.
A problemática da Dengue, tanto endêmica quanto epidêmica é aderente às Diretrizes Curriculares Nacionais no que se refere às especificações IV, V, VI, VIII, XVI, XX, XXIX das competências específicas.
A questão exigia um nível elevado de profundidade evidenciando/valorizando o quadro clínico com ênfase na fisiopatologia.
Acredita -se que isso levou a interpretações superficiais e inadequadas de resposta ao item ‘ a ’.
A questão foi dividida em dois subitens : o primeiro avaliando o conhecimento inerente à fundamentação científica que sustenta a justificativa da sintomatologia apresentada.
O segundo item descrevendo a assistência de acordo com o protocolo do Ministério da Saúde.
Isso contribuiu tanto na avaliação do conteúdo e adequação das respostas, considerando a complexidade do caso, quanto para a pontuação/totalização das notas.
Muitos estudantes, sem se atentar à informação do enunciado : “ diagnóstico confirmado de dengue hemorrágica ” traçaram o plano de cuidados como se o paciente estive em avaliação ambulatorial, com sintomas de Dengue Clássica, recomendando cuidados domiciliares quanto à hidratação, analgesia adequada e sinais de alarme, com necessidade de retorno precoce.
Como a pergunta do item ‘ a ’ contém um “ e ”, esperávamos que o aluno se pautasse tanto no quadro clínico quanto nos resultados dos exames.
No entanto, o padrão de resposta não trazia a exigência dessa associação, já que elencava oito grupos de sintomas que estão relacionados ao quadro clínico ou aos exames laboratoriais. 1.
A febre como sintoma inicial foi resultante do processo de infecção viral aguda e resposta imunológica à presença de patógeno. 2.
Os quadros de dor muscular, articular e retro-orbitária, são resultantes do processo de febre alta que gera letargia e fadiga muscular, além de caracterizarem o quadro clínico da dengue. 3.
A dor abdominal é causada por hepatomegalia ; 4.
A gengivorragia é resultante da plaquetopenia e fragilidade capilar dos pequenos vasos sanguíneos \textbf{desencadeada} pela fricção durante a escovação. 5.
O quadro de náuseas e vômitos é resultante da dor, febre e compressão gástrica causada pela hepatomegalia. 6.
A letargia é decorrente do quadro infeccioso geral e pela perda hidroeletrolítica gerada pelo vômito. 7.
Os exames laboratoriais apresentaram anemia e plaquetopenia, decorrentes da diminuição da função hepática. 8.
A prova do laço evidenciou resultado expressivo pela plaquetopenia e fragilidade capilar decorrente das alterações vasculares causados pelo quadro clínico da doença.
Os itens 4 e 8 do padrão se referem a sintomas decorrentes da plaquetopenia e foram pontuados separadamente, já que os demais eram mais associados ao quadro clínico.
Na subdivisão da pontuação, alguns itens presentes nas respostas dos alunos valeram poucos pontos apesar de apresentarem vários sintomas como ocorre no item 2.
Ademais, a letargia poderia também ter sido mais bem pontuada quando o estudante apontava a baixa oxigenação decorrente do declínio da taxa dos eritrócitos como razão para a mesma.
Alguns estudantes entenderam que no item ‘ a ’ ( justificar a sintomatologia com base no quadro clínico ) deveriam justificar que a paciente tinha Dengue devido aos sintomas apresentados na introdução da questão discursiva, a saber : febre, hepatomegalia, mialgia, artralgia etc.
Dessa forma, não justificavam essas sintomatologias, mas apenas afirmavam que esses são os sintomas da dengue.
O padrão de resposta do item ‘ b ’, com pergunta bastante apropriada para o objeto de trabalho do enfermeiro, também elencava oito itens : 1.
Realizar notificação compulsória ; 2.
Estabelecer acesso venoso periférico calibroso para reposição hidroeletrolítica e manter acesso vascular para administração de medicamentos, conforme protocolo do MS ; 3.
Realizar controle rigoroso de temperatura, além do uso de compressas mornas e promoção da ingesta hídrica para reposição de volume no caso de febre. 4.
Instalar oxigenoterapia para compensação da capacidade de oxigenação tecidual reduzida pela eritrocitopenia e hemoglobinopenia ; 5.
Controlar o balanço hidroeletrolítico, com atenção para diurese em relação a quantidade, aspecto e coloração que poderão indicar hematúria como sinal de complicação do quadro de dengue hemorrágica ou oligúria como evolução de choque ; 6.
Avaliar locais de punções para sinais de sangramentos e evitar procedimentos invasivos por razão da plaquetopenia, bem como cuidados com a gengiva pra evitar novos episódios sangrantes ; 7.
Garantir avaliação periódica da realização dos exames laboratoriais para acompanhar evolução da competência plaquetária e reestabelecimento dos valores hemograma ; 8.
Avaliar os sinais indicativos de evolução de quadro de choque que poderão instalar por ocasião de quadro progressivo de dengue hemorrágica.
No entanto, algumas respostas possíveis não foram contempladas pelo padrão ( e, portanto, não receberam pontuação ) : a internação ( paciente estava no PS ) ; encaminhar para o repouso ( devido ao risco de hemorragia e a necessidade de preservação das reservas energéticas de um organismo que se encontra desgastado pelo vírus ) ; avaliação do ambiente domiciliar e de trabalho na busca de focos do mosquito ( busca ativa ) ; assistência referente a dor, náusea, vômito.
A descrição da condição clínica é objetiva e o quadro sintomatológico associado aos valores laboratoriais fornecia, para o estudante, informações necessárias para desenvolver a resposta.
Entretanto, os resultados do hemograma não foram analisados corretamente, em concomitância com a sintomatologia descrita no enunciado.
Os valores de alguns componentes ( como eritrócitos, hemoglobina, linfócitos atípicos ) não foram bem interpretados.
O mesmo ocorreu para a não distinção entre leucopenia e leucocitose e inflamação e infecção.
A assistência de enfermagem, nesse caso, tem por objetivo identificar precocemente as formas graves da doença, orientar o plano de acompanhamento e a conduta terapêutica adequada para cada situação clínica, propondo uma abordagem clínica evolutiva.
O acolhimento e a identificação de risco social e o histórico epidemiológico não foram citados pelos alunos, bem como a exigência de Notificação Compulsória.
O padrão de resposta não contemplava plenamente as fases de tratamento, prevenção e monitoramento de sinais de choque para a fase da doença em que a cliente do caso se encontrava ( fase C de alarme – sinal hemorrágico ).
Além disso, não previa outras condutas apontadas pelos estudantes como importantes para a assistência de enfermagem que poderiam ser consideradas corretas.
Por exemplo, o enunciado informa que a paciente faz uso de dipirona, entretanto, a ação de enfermagem de suspender o medicamento ou orientar a não administração do mesmo não fazia parte do padrão.
Outras respostas não contempladas envolviam : aconselhar que a paciente ficasse isolada para não contaminar outras pessoas ; comunicação ao MS ; recomendação de fazer a varredura para saber onde está o foco ( residência/trabalho ) que contaminou a paciente ; recomendação de repouso da paciente.
Destaca -se um grande número de respostas contendo apenas uma assistência/prescrição/intervenção de enfermagem que foi “ encaminhamento ao médico ”, o que não contempla uma ação de enfermagem.
Foram encontradas explicações totalmente equivocadas e distorcidas de terminologias que até mesmo estudantes de períodos mais precoces já deveriam ter apreendido.
Também houve respostas com prescrições de enfermagem inaceitáveis, as quais resultariam em danos à saúde da paciente, tais como : hepatomegalia definida como inflamação do fígado ; plaquetopenia conceituada como se todas as células sanguíneas estivessem alteradas ; prescrição de realizar exercício físico e deambulação ; orientações do tipo eliminar reservatórios de água parada e depois encaminhar a paciente para residência ; plaquetas baixas devido ao sangramento ( e não o contrário ).
Muitas respostas eram vagas, como por exemplo : acompanhar o quadro clínico, manter em \textbf{observação} ; checar mudança do quadro/sinais e sintomas ; fazer exame físico/anamnese ; cuidados gerais ; dieta adequada/saudável ; uso de EPI.
Os estudantes, de uma maneira geral, evidenciaram dificuldade para responder a Questão 4.
Dentre os que responderam, grande parte optou por um sintoma ou dois dos listados no caso clínico e justificou com base na fisiopatologia.
A maior dificuldade foi apresentar justificativa para a sintomatologia.
Um número expressivo de estudantes optou por listar novamente os problemas apresentados no enunciado, sem fundamentação ou apresentando justificativas incorretas.
Na prescrição de ações que caracterizam a assistência de enfermagem, em atendimento ao Protocolo do Ministério da Saúde, as respostas foram vagas, relacionadas com um cuidado padrão e ou rotineiro, e não específicas para o caso.
As sintomatologias mais descritas foram a hemorragia ( baixa de plaqueta ) e a infecção que leva à febre – leucopenia.
Como consequência, os cuidados mais prescritos foram hidratar controle da hipertermia.
A Questão 4 apresentou de forma objetiva, situações que requerem uma intervenção imediata e segura demandando do estudante a correta identificação dos problemas e riscos, o que não aconteceu. 3.3.2.5 Análise da Questão Discursiva 5 do Componente de Conhecimento Específico A Tabela 3.15 contém as informações relativas à questão 5 do conjunto do Componente de Conhecimento Específico.
O desempenho dos estudantes nessa questão foi inferior ao da questão 3 e superior ao da questão 4.
A nota média dos estudantes de todo o Brasil foi 31,8.
A maior média foi registrada na região Sul ( 34,5 ), enquanto a menor média foi registrada na região Norte ( 29,6 ).
Quanto à variabilidade das notas, o desvio padrão dos alunos do Brasil, como um todo, foi 30,0.
Enquanto o maior desvio foi encontrado na região Nordeste ( 30,9 ), o menor foi encontrado na região Sudeste ( 29,2 ).
A mediana ( 30,0 ), a nota máxima ( 100,0 ) e a nota mínima ( 0,0 ), foram as mesmas para todas as regiões do Brasil, exceto a mediana na região Sul ( 35,0 ).
Os estudantes, de uma maneira geral, evidenciaram dificuldade para responder a Questão 4.
Dentre os que responderam, grande parte optou por um sintoma ou dois dos listados no caso clínico e justificou com base na fisiopatologia.
A maior dificuldade foi apresentar justificativa para a sintomatologia.
Um número expressivo de estudantes optou por listar novamente os problemas apresentados no enunciado, sem fundamentação ou apresentando justificativas incorretas.
Na prescrição de ações que caracterizam a assistência de enfermagem, em atendimento ao Protocolo do Ministério da Saúde, as respostas foram vagas, relacionadas com um cuidado padrão e ou rotineiro, e não específicas para o caso.
As sintomatologias mais descritas foram a hemorragia ( baixa de plaqueta ) e a infecção que leva à febre – leucopenia.
Como consequência, os cuidados mais prescritos foram hidratar controle da hipertermia.
A Questão 4 apresentou de forma objetiva, situações que requerem uma intervenção imediata e segura demandando do estudante a correta identificação dos problemas e riscos, o que não aconteceu. 3.3.2.5 Análise da Questão Discursiva 5 do Componente de Conhecimento Específico A Tabela 3.15 contém as informações relativas à questão 5 do conjunto do Componente de Conhecimento Específico.
O desempenho dos estudantes nessa questão foi inferior ao da questão 3 e superior ao da questão 4.
A nota média dos estudantes de todo o Brasil foi 31,8.
A maior média foi registrada na região Sul ( 34,5 ), enquanto a menor média foi registrada na região Norte ( 29,6 ).
Quanto à variabilidade das notas, o desvio padrão dos alunos do Brasil, como um todo, foi 30,0.
Enquanto o maior desvio foi encontrado na região Nordeste ( 30,9 ), o menor foi encontrado na região Sudeste ( 29,2 ).
A mediana ( 30,0 ), a nota máxima ( 100,0 ) e a nota mínima ( 0,0 ), foram as mesmas para todas as regiões do Brasil, exceto a mediana na região Sul ( 35,0 ).
O Gráfico 3.22 representa a distribuição das notas da questão discursiva 5 do Componente de Conhecimento Específico.
Destaca -se, como nas demais questões discursivas, o grande número de estudantes que deixaram esta questão em branco, correspondendo à moda da distribuição.
Entre as três questões discursivas do Componente de Conhecimento Específico, esta foi a que apresentou um maior espalhamento das notas daqueles que resolveram a questão.
Além da moda na nota zero, aparecem ainda dois máximos locais : ( 30 ; 40 ] e ( 60 ; 70 ].
O coeficiente de assimetria para todos os alunos é 0,29 e para todas as regiões se mantém positivo, variando entre 0,14 na região Sul e 0,44 na região Norte. 3.3.2.6 Comentários sobre as respostas à Questão Discursiva 5 A Questão 5 era clara e continha informações suficientes e adequadas para promover a reflexão e a crítica dos estudantes de enfermagem acerca das demandas de gestão e de organização do Sistema Único de Saúde.
Com linguagem pertinente, o enunciado apresentava os fatos e os dados, explorando conteúdos consonantes com as Diretrizes Curriculares Nacionais do Curso de Enfermagem.
O nível de dificuldade foi considerado mediano no que tange ao conteúdo teórico e às experiências teórico-práticas e às práticas da graduação em Enfermagem.
Entretanto, as respostas dos alunos evidenciaram inexpressiva compreensão do conceito de territorialização e de diagnóstico situacional, bem como sua articulação com um plano de intervenção voltado para resolver os problemas de saúde apontados no enunciado.
O item ‘ a ’ solicitava que se enumerassem os passos a serem seguidos para a realização da territorialização e do diagnóstico situacional do município.
A enumeração indica uma ordenação de passos que não foi seguida pelos respondentes.
Mesmo com um padrão de resposta que possibilitava pontuar parte dos passos necessários e os ajustes necessários de adequação às respostas dos estudantes, o que foi respondido era sempre insuficiente.
Detectou -se muita dificuldade para responder ao item ‘ a ’, evidenciando desconhecimentos acerca do conceito e da realização do processo de territorialização e, em decorrência, também do diagnóstico situacional do município.
No item ‘ b ’ foi solicitado escolher dois problemas, dentre os que já estavam relacionados no enunciado da questão, e traçar um plano de intervenção para resolvê -los.
O padrão de resposta era bastante flexível de modo que foi possível atribuir alguma pontuação para a maioria das respostas oferecidas pelos estudantes.
Em todas as respostas, de algum modo, foram abordados/citados problemas/intervenções com predominância da taxa de mortalidade materna, taxa de mortalidade infantil e falta de registro de casos de hanseníase e tuberculose.
Apesar disso, houve um grande número de respostas erradas.
Para o item ‘ a ’, as respostas encontradas com maior frequência foram as que citavam cadastramento de pessoas da comunidade, habitantes do município ; levantamento de número de pessoas, doentes, crianças, idosos ; identificação de serviços ; aplicação de formulários e realização de entrevistas, atendendo ao passo um da territorialização e diagnóstico situacional.
Houve ainda, em menor proporção, o uso das expressões como “ áreas de risco ” e “ mapeamento ” que indicaram o alcance do passo dois.
Em relação ao item ‘ b ’, foi maciça a escolha de problemas referentes à alta taxa de mortalidade infantil, alta taxa de mortalidade materna e falta de registro de casos de hanseníase e tuberculose.
Cabe observar que havia interseção entre a Questão Discursiva 5 e algumas das questões objetivas, a saber : Questão 29 ( alternativas C e D ), Questão 33 ( porque aborda a problemática do controle e tratamento da tuberculose ) e Questão 34 ( com destaque para as alternativas B e D ).
No entanto, concluída a correção, observou -se que foram raras as referências à Rede Cegonha como espaço de intervenção para a redução da Mortalidade Materna e Perinatal.

% matched lemmas: aceleração, apenar, chão, câmera, desencadear, disputar, observação, peso
\end{textsample}
